% Pursuing Stacks (cotes 134-2 to 134-8): archive cote 134-4, site page 54 (PDF page 55), typescript page 142.
% typescript. Draft compared: ps3.tex:2730-2802.
% AI-transcribed by Claude Opus 5.5 (claude-opus-5-5, adaptive thinking, effort high), Message Batches API, batch ps2 (msgbatch_01NVwr1Yuo8qJf5rykeGBXPK), 2026-11-12; not yet checked by eye.
\scanpage{134-4/54}
\mspage{142}

\uline{Proposition} Let M be a category \deleted{with finite products}, \deleted{and} \inserted{assume M,} \uline{totally}
\uline{O-connected} (i.e. satisfying conditions \inserted{a)}b)c)e) of page 123). Let W$\hand{\subset}$F$\hand{\ell}$(M)
be a saturated set of arrows in M, consider the associated homotopy
structure $h_W$ on M, with W-aspheric intervals as a generating family of
weak homotopy intervals. Consider also the canonical homotopy structure
\xout{xxxM} $h_M$ on M, with O-connected intervals as a generating family of weak
homotopy intervals (thus $h_M = h_{W_o}$ , where $W_o \hand{\subset}$ F$\hand{\ell}$(M) consists of all
arrows made invertible by the functor $\hand{\pi}_o$ : M$\hand{\longrightarrow}$(sets)). In order for the
equality $h_W = h_M$ to hold, it is necessary and sufficient that the
following two conditions be satisfied:

a) Any object \xout{with} I of M which is W-aspheric over e and admits a
section, is connected (it is enough for this that W be O-admissible i.e.
f $\hand{\in}$ W imply $\hand{\pi}_o$(f) bijective);

b) For any connected object I of M and any two sections $\hand{\delta}_o,\hand{\delta}_1$ of I,
these can be \hand{``}joined\hand{''} by \deleted{a} \inserted{finite} chain of sections $s_i$ ( O$\hand{\leq}$ i$\hand{\leq}$ n) , $s_i = \hand{\delta}_o$,
$s_n = \hand{\delta}_1$ , such that for any two consecutive ones , there exists an obje\hand{ct}
J \xout{\unsure{of M asph}} W-aspheric over e, a map J$\hand{\longrightarrow}$I and two section of J mapped
into the sections $s_i,s_{i+1}$ of I.

THe condition a) says there are not too many weak equivalences, \xout{wher}
whereas b) says there are still enough for "testing" connectedness in
terms of \over{\unsure{W}}{W}-aspheric intervals. Both conditions look plausible enough !

Next step one would think of, in this context, is to give conditions
on W (independently of the previous ones) which will allow to express
\margin{(*) Compare with p. 135.}
W in terms of $h_W$ \deleted{\illegible{1}}. But for this, I should develop first a description
of a notion of "weak equivalence" in terms of an arbitrary homotopy struc-%
ture h on M, as contemplated in the previous paragraph\hand{(*)}. I decided howe-%
ver to give precedence to the "key result" still ahead, \xout{\unsure{at any rate}}
\xout{\unsure{Certain}} Maybe the condition $h_W = h_M$ , expressed in the previous proposi-%
tion, merits a name - we'll say that the notion of weak equivalence W
is "\uline{geometric}", if it satisfies the two conditions a) and b) above,
or rather the slightly stronger a') in place of a) - namely O-admissibi-%
lity of W (plus b) of course). T\over{\unsure{h}}{\unsure{h}}e conditions a\hand{'}) and b) are the expli-%
cit ones for checking - but for using that W is geometric, the
more conceptual statement\xout{s} $\over{h}{h}_W = h_M$ \inserted[typed]{(besides O-admissibility)} is the best. Thus, as any \uline{contrac-}%
\margin{$+$}
\uline{tor} in M which is connected (hence a weak homotopy interval for $h_M$) is
$h_M$-contractible, it is $h_W$-contractible, and hence W-aspheric over e,
and co\over{\unsure{w}}{n}versely of course. Using this, we get a converse to part c) and

\margin{54}
